\documentclass[12pt,fleqn,a4paper]{article}

\usepackage{latexsym}
\usepackage{url}
\usepackage{xspace}
\usepackage{epsfig}
\usepackage{psfrag}
\usepackage{a4wide}
\usepackage{marvosym}
\usepackage{amsmath,amsfonts,amssymb,amsthm,latexsym}
\usepackage{graphics,graphicx,color,subfigure}
\usepackage{fancyhdr}
\usepackage[english]{babel}
\usepackage[latin1]{inputenc}

\textheight 680pt
\textwidth 460pt
\topmargin -40pt
\oddsidemargin 5pt
\evensidemargin 5pt
\parindent 0pt

\pagestyle{fancyplain} \setlength{\headheight}{16pt}
\renewcommand{\sectionmark}[1]{\markright{\thesection\ #1}}
\lhead[\fancyplain{}{\thepage}]
    {\fancyplain{}{\rightmark}}
\rhead[\fancyplain{}{\leftmark}]
    {\fancyplain{}{\thepage}}
\cfoot{}
\renewcommand{\thesection}{\arabic{section}}
\renewcommand{\thesubsection}{\arabic{section}.\arabic{subsection}}


\begin{document}
\begin{titlepage}%Institution
\vspace{2cm}
\centerline{
\large{Department of Artificial Intelligence and Human Interfaces}}
\vspace{0.2cm}
\centerline{\large{University of Salzburg}}%Title with one or two Lines(More if wanted)
%\hline
\vspace{2cm}

\centerline{\large{PS Natural Computation}}
\centerline{SS 26}
\vspace{1cm}

\centerline{\Large{\bf{AF Basics Group}} }%Type of the Document
\vspace{1cm}

\vspace{0.4cm}%Date
\centerline{\today}
\vspace{5cm}%Authors

%\hline
\vspace{0.2cm}
Project Members:\\
\centerline{Gerwin Bacher}\\
\centerline{Andreas Schlager}\\
\centerline{Linda Schmidt}\\
\centerline{Tomas Tomandl}\\
\centerline{Sophie Wagner}\\
\vspace {1cm}\\

Academic Supervisor: \\
\centerline{Helmut MAYER}
\centerline{helmut@cosy.sbg.ac.at}
\vspace{1.5cm}\\
Correspondence to: \\
\centerline{Universit\"{a}t Salzburg} \\
\centerline{Fachbereich AIHI} \\
\centerline{Jakob--Haringer--Stra\ss e 2} \\
\centerline{A--5020 Salzburg} \\
\centerline{Austria}
\clearpage
\end{titlepage}

%Table of Content
% \setcounter{page}{1}
% \pagenumbering{Roman} %I,II,III... in the TOC
% \tableofcontents

\clearpage
\pagestyle{headings}
\pagenumbering{arabic}  %Better if TOC is variable (more than one page)
\setcounter{page}{1}
\pagenumbering{arabic}  %Better if TOC is variable (more than one page)
\setcounter{page}{1}

\abstract{
Artificial neural networks are widely used in modern machine learning for classification 
and pattern recognition tasks. Their performance depends on several design choices, 
such as the network architecture and the activation functions used in the neurons. This 
project examines how different activation functions affect neural network performance 
on classification tasks of varying difficulty. 
Three benchmark problems were defined for all participating groups. These include the 
classification of concentric circles as an easy task , the two-moons problem as an 
medium hard task and overlapping spirals as a more challenging one. To establish a 
common baseline for later experiments, our group evaluates six activation functions: 
ReLU, Sine, GeLU, Sigmoid, TanH, and Swish, using identical network architectures. In 
each case, a single activation function is applied consistently across all neurons in the 
network. 
The results of this baseline study will serve as a reference for other groups investigating 
whether adapting activation functions at the neuron level can improve overall neural 
network performance.
}

\section{Introduction}
Natural computation studies computational methods inspired by processes found in 
nature. A central example of this field is artificial neural networks, which are modeled 
after the structure and functioning of biological brains. These networks learn patterns 
from data and are widely used in classification, prediction, and optimization tasks. 
A neural network consists of interconnected neurons arranged in layers. Each neuron 
processes incoming signals and passes its output to the next layer. A key component in 
this process is the activation function, which determines how strongly a neuron responds 
to its input and introduces nonlinearity into the model. Without such nonlinearities, 
neural networks would be limited to solving only linearly separable problems. 
Different activation functions exhibit different mathematical properties, which in turn 
influence how a network learns. Some lead to faster convergence, while others improve 
stability or enable the learning of more complex patterns. In this work, we compare six 
commonly used activation functions: ReLU, Sine, GeLU, Sigmoid, TanH, and Swish. 
To analyze their behavior, three classification problems with increasing difficulty were 
defined. The first consists of two concentric circles and represents a relatively simple 
nonlinear task. The second is the well-known two-moons problem, which typically 
achieves an accuracy of around 85 to 95 percent. The third problem involves two 
overlapping spirals, a benchmark known for its high level of nonlinearity and difficulty. 

To ensure fair comparisons, the neural network architecture was fixed across all 
experiments and participating groups, including the number of layers and neurons per 
layer. This allows the effects of different activation functions to be observed more directly. 
The main goal of this project is to investigate whether neural networks can achieve better 
performance when activation functions are adapted individually for each neuron rather 
than being fixed across the entire model. As a first step, however, a solid baseline is 
required. Therefore, this study focuses on training and evaluating networks that use a 
single activation function consistently. 
The results provide a foundation for future comparisons with adaptive activation 
approaches and may help identify which activation functions are better suited to specific 
types of classification problems.
\newpage

\section{Milestones} % NOT in final paper
\begin{itemize}
\item{Due date:}
\end{itemize}

\newpage

\section{Progress of Work} % NOT in final paper

\subsection{Week 1, Tuesday, 18.10.2005}

\subsection{Week 2, Tuesday, 25.10.2005}

\subsection{Week 3, Tuesday, 08.11.2004}

\subsection{Week 4, Tuesday, 22.11.2005}

\subsection{Week 5, Tuesday, 29.11.2005}

\subsection{Week 6, Tuesday, 06.12.2005}

\subsection{Week 7, Tuesday, 13.12.2005}

\subsection{Week 8, Tuesday, 20.12.2005}

\subsection{Week 9, Tuesday, 10.01.2006}

\subsection{Week 10, Tuesday, 17.01.2006}

\subsection{Week 11, Tuesday, 24.01.2006}

\newpage

\section{Subproject Responsibilities} % NOT in final paper
\newpage

% links go here, NOT in references
\section{Links}

\begin{itemize}
\item Project Page: \url{http://student.cosy.sbg.ac.at/???}
\item PS Page:
\url{http://www.cosy.sbg.ac.at/~helmut/Teaching/PSRobotik/}

\end{itemize}

% real scientific references
\bibliography{}		% .bib files here

\end{document}
